\chapter{Experimental Setup}
\section{Hardware and Measurement Conventions}
The implementation targets CUDA-capable GPUs and is specifically optimized for SM90/H100 systems in the stable tensor-core path. Two measurement regimes are used:
\begin{enumerate}[leftmargin=2em]
\item \texttt{BSI\_PROFILE=1} for build/dot attribution,
\item \texttt{BSI\_PROFILE=0} for fairer runtime timing.
\end{enumerate}
Because dot time is the primary optimization target in this thesis, the main kernel-attribution results use \texttt{BSI\_PROFILE=1}. Those numbers should be interpreted as internal timing breakdowns rather than fully asynchronous throughput measurements.

\section{Primary Evaluation Configuration}
Unless otherwise noted, the principal results use:
\begin{itemize}[leftmargin=2em]
\item \texttt{decimal\_places = 2}
\item \texttt{compress\_threshold = 0.5}
\item \texttt{num\_samples = 1000}
\item \texttt{max\_seq\_len = 512}
\item \texttt{base\_dtype = fp16}
\end{itemize}

\section{Models}
The primary model family is OPT with the following checkpoints:
\begin{itemize}[leftmargin=2em]
\item \texttt{facebook/opt-125m}
\item \texttt{facebook/opt-1.3b}
\item \texttt{facebook/opt-6.7b}
\end{itemize}
The 30B model is used as a stress test to study scope sensitivity rather than as the primary optimization target.

\section{Representative Kernel Shapes}
Kernel-only studies emphasize two shapes for the stable baseline:
\begin{itemize}[leftmargin=2em]
\item $Q=512, R=16384, D=4096$, corresponding to an FC1-like layer in \texttt{opt-6.7b}
\item $Q=512, R=50272, D=4096$, corresponding to a very wide output or vocabulary-scale projection
\end{itemize}
These shapes are also used as the primary Nsight Compute profiling targets.

\section{Scope Policies}
Three scope policies are implemented:
\begin{itemize}[leftmargin=2em]
\item \texttt{scope=all}: replace all eligible linear layers,
\item \texttt{scope=attention}: replace only attention projections,
\item \texttt{scope=mlp}: replace only feed-forward layers.
\end{itemize}
The accuracy results show that these scope choices correspond to meaningfully different performance and quality regimes.

\section{Documentation and Reproducibility Workflow}
The operational workflow used throughout the study consists of four steps:
\begin{enumerate}[leftmargin=2em]
\item initialize network and cluster-specific GPU environment helpers,
\item activate the Python environment and rebuild the extension,
\item run kernel-only and model-level experiments,
\item profile the stable kernel with Nsight Compute.
\end{enumerate}
The exact commands are collected in Appendix~A and are mirrored in the project README so that the operational documentation stays aligned across the written report and the project documentation.
